\section{Quantum Information, Science \& Technology}
\subsection[Challenge 7: Noisy quantum sensing]{Challenge 7: Noisy quantum sensing}
\textbf{Problem ID:} \texttt{Challenge\_7\_main}\\
\textbf{Primary inferred discipline:} Quantum Information, Science \& Technology\\
\textbf{Secondary inferred disciplines:} Atomic, Molecular \& Optical\\
\textbf{Python implementation:} \texttt{challenges/07/answer.py}

\subsubsection*{Problem}
\paragraph*{Problem setup}

Consider a noisy distributed quantum sensing scenario as follows. The quantum sensor network contains $d$ sensor nodes, and each node contains $n$ sensor qubits. There is a potentially different parameter to be encoded on each sensor node locally. Our sensing objective is to estimate the scaled average of all $d$ local parameters, i.e., supposing that the parameters are $x_1, x_2, \dots, x_d$, we want to estimate $\theta_1 = (\sum_{i=1}^d x_i) / \sqrt{d}$. The quantum network will distribute a $d$-qubit Greenberger-Horne-Zeilinger (GHZ) state to the $d$ sensor nodes, while each sensor node will hold one qubit of the GHZ state. Then each sensor node further performs local entanglement generation between the qubit from the $d$-qubit GHZ state and the remaining $n-1$ sensor qubits. The final state is the initial probe state, which we assume to be an $nd$-qubit noisy GHZ state across the quantum sensor network. Suppose that the global GHZ state is in a depolarized form, with GHZ fidelity $F(n)=Fk^{n-1}$. Now assume that the noisy sensing dynamics can be described by independent single-qubit Lindblad equations for all $nd$ sensor qubits: $\frac{d}{dt}\rho = -i\frac{\omega^{(i)}}{2}\left[\sigma_z^{(i,k)},\rho\right] + \frac{\gamma}{2}\left(\sigma_z^{(i,k)}\rho \sigma_z^{(i,k)} - \rho\right)$ for the $k$-th qubit on the $i$-th node, where $\omega^{(i)}$ is the precession frequency of each qubit sensor on node $i$, and $\gamma$ is the single-qubit dephasing rate. In other words, we assume that the dephasing rate is homogeneous across the entire sensor network, while the angular frequency is the same for each qubit on one node but can be generally different between nodes.

\paragraph*{Main problem}

Suppose that the noisy sensing dynamics has duration $t$, and each local parameter is the accumulated phase, i.e. $x_i = \omega^{(i)} t$. What is the quantum Fisher information for $\theta_1$? Do not explicitly include $\gamma$ and $t$ in the final expression. Instead, you may use another variable $q = (1 + e^{-\gamma t}) / 2$.

{
\color{blue}
\paragraph*{Issues with the problem statement}

The main problem is the sentence \emph{"Suppose that the global GHZ state is in a depolarized form, with GHZ fidelity $F(n)=Fk^{n-1}$"} which is ambiguous for the following reasons:
\begin{itemize}
  \item The sentence tells the reader/AI to assume the initial state to be a GHZ state that went through a depolarizing channel. This in general means a density matrix of the form
  \begin{equation} \label{eq:depol_state}
      \rho\, = (1-p)\,\ket{\mathrm{GHZ}_N}\!\!\bra{\mathrm{GHZ}_N} + p \frac{1}{2^N} I
  \end{equation}
  where $p$ is the depolarizing probability, $N=n\,d$ the number of qubits and $I$ the identity operator. If we are interested in the fidelity with the pure GHZ state then it follows the function ${F(N)=1-p\,\big(1-2^{-N}\big)}$. If the instruction is instead assumes a tensor product of single-qubit depolarizing channels (which is not equivalent to the global one just described), the fidelity will still not be described by the provided function.

  \item Even if one assumes the given fidelity function to be true, is not clear what the objects $F$ and $k$ are. Are they scalars? Are they functions of $n$ or $d$?  
\end{itemize}
This ambiguity does not make the problem is unsolvable in general. It just adds confusion and unnecessary information, which I believe could lead an AI agent in some rabbit hole (which is what I believe happened in this case).

Another minor thing about the problem statement: the last two sentences stating the requirement about $\gamma$, $t$ and $q$ may be constraining for an AI agent as it will blindly try to stick to this notation from the get go. In contrast a human solver may ignore it until the very end and apply the substitution only for the final answer (this is what I do in my solution for example).

\paragraph*{How I would have solved the problem}

I would solve the problem in the following steps:
\begin{enumerate}
  \item Note the ambiguity and decide to treat the problem while assuming that the initial state $\rho$ is given by Eq.~\eqref{eq:depol_state}.
  
  \item Note that the dynamics of the sensor commutes with the depolarizing channel. So one can first focus on an initially pure state $\ket{\mathrm{GHZ}_N}$ instead of the depolarized one.
  
  \item Note that the sensing Hamiltonian and the dephasing dynamics commute, thus we can further split it into a unitary part followed by a decoherence channel. In other words the dynamics of the state can be expressed as
  \begin{equation}
    \rho_\mathrm{end} = \mathcal{D}_p\circ\mathbf{\Phi}_t\circ\mathcal{U}_t \left( \ket{\mathrm{GHZ}_N}\!\!\bra{\mathrm{GHZ}_N} \right)
  \end{equation}
  where $\mathcal{U}_t$ is the propagator of the sensing Hamiltonian up to time $t$, $\mathbf{\Phi}_t$ is the dephasing dynamics integrated up to time $t$ and $\mathcal{D}_p$ is the initial depolarizing channel with probability $p$.

  \item Observe that $\mathcal{U}_t$ just adds a relative phase (parametrized by $\theta_1$) to the GHZ state, i.e.
  \begin{equation} \label{eq:deph_ghz_phase}
    \ket{0}^{\otimes N}+\ket{1}^{\otimes N} \longmapsto
    \ket{0}^{\otimes N}+e^{i \varphi}\ket{1}^{\otimes N}
    \qquad\text{with}\qquad \varphi=n\sqrt{d}\,\theta_1
  \end{equation}

  \item So the problem boils down to a standard result in quantum sensing which we can directly apply here (see for example W. Zong et al. Phys. Rev. A 87, 022337 (2013)):
  
  The Quantum Fisher Information for a GHZ state of size $N$ with a phase $\phi$ evolving though a dephasing channel $\mathbf{\Phi}_t$ at a rate $\gamma$ is (Eq.~(87) in previous Ref.) 
  $$
    \mathcal{F}_\varphi = (\partial_{\theta_1}\varphi) \ e^{-2\,N^2\,\gamma\,t} = n^2\,d\,(2q-1)^{2nd}
  $$
  Adding in the initial depolarizing channel (which we can do by multiplying the transformation matrices $A$ for $\mathcal{D}_p$ and $\mathbf{\Phi}_t$), we would get
  \begin{equation} \label{eq:my_final_answer}
    \mathcal{F}_\varphi = n^2\,d\,(2q-1)^{2nd}\frac{(1-p)^2}{\frac{2}{d} + (1-p)\,G(\omega)}  
  \end{equation}
  where the second term is derived using  Eq.~(60-62) and where $G(\omega)$ has to be further simplified for a GHZ state (I leave it as is since it is of little importance). I will discuss this term in my review of the AI response below.

  Equation~\eqref{eq:my_final_answer} is my final answer to the problem statement.

\end{enumerate}

\textbf{Disclosure:} I have \textbf{not} used any AI model to brainstorm my approach or derive Eq.~\eqref{eq:my_final_answer}, and have not looked at the solution prior to formulate my answer. 


}